\documentclass[10pt,a4paper]{amsart}
\usepackage{amsmath,amssymb,amsthm,graphicx}
\RequirePackage{amssymb,amsmath,amsthm,bm}
\RequirePackage{thmtools}
\RequirePackage[capitalize]{cleveref}
\providecommand{\tightlist}{%
  \setlength{\itemsep}{0pt}\setlength{\parskip}{0pt}}

%%%%% varphi instead of phi
\renewcommand{\phi}{\varphi}

\mathchardef\mhyph="2D 
\providecommand{\A}{\mathcal{A}}
\providecommand{\C}{\mathcal{C}}
\providecommand{\L}{\mathcal{L}}
\providecommand{\B}{\mathcal{B}}
\providecommand{\M}{\mathcal{M}}
\renewcommand{\S}{\mathcal{S}}
\providecommand{\ol}[1]{\overline{#1}}
\providecommand{\mbf}[1]{\mathbf{#1}}
\providecommand{\mf}[1]{\mathfrak{#1}}
\providecommand{\mrm}[1]{\mathrm{#1}}
\providecommand{\mc}[1]{\mathcal{#1}}
\providecommand{\ra}{\rightarrow}
\providecommand{\Ra}{\Rightarrow}
\providecommand{\La}{\Leftarrow}
\providecommand{\lr}{\leftrightarrow}
\providecommand{\hf}{\mathbb{HF}}
\providecommand{\spec}{DgSp}
\providecommand{\bSigma}{\bm\Sigma}
\providecommand{\bDelta}{\bm\Delta}
\providecommand{\bPi}{\bm\Pi}
\providecommand{\Pinf}[1]{\Pi^{\mathrm{in}}_{#1}}
\providecommand{\Sinf}[1]{\Sigma^{\mathrm{in}}_{#1}}
\providecommand{\dSinf}[1]{d\mhyph\Sigma^{\mathrm{in}}_{#1}}
\providecommand{\Dinf}[1]{\Delta^{\mathrm{in}}_{#1}}
\providecommand{\bp}{\bar p}
\providecommand{\bq}{\bar q}
\providecommand{\ba}{\bar a}
\providecommand{\bb}{\bar b}
\providecommand{\bc}{\bar c}
\providecommand{\bd}{\bar d}
\providecommand{\bx}{\bar x}
\providecommand{\by}{\bar y}
\providecommand{\bz}{\bar z}
\providecommand{\LO}{\mathrm{LO}}
\declaretheorem[name={Theorem},numberwithin=section]{theorem}
\declaretheorem[name={Proposition}, sibling=theorem]{proposition}
\declaretheorem[name={Lemma}, sibling=theorem]{lemma}
\declaretheorem[name={Fact}, sibling=theorem]{fact}
\declaretheorem[name={Claim}, numberwithin=theorem]{claim}
\declaretheorem[name={Definition},numberwithin=section,sibling=theorem, style=definition]{definition}
\declaretheorem[name={Remark}, numbered=no,style=definition]{remark}
\declaretheorem[name={Corollary}, sibling=theorem]{corollary}
\declaretheorem[name={Question},sibling=theorem]{question}
\newcommand{\T}{{\mathcal{T}}}
\newcommand{\SR}{SR}
\newcommand{\pSR}{pSR}
\newcommand{\Q}{\mathbb{Q}}
\newcommand{\Z}{\mathbb{Z}}
\newcommand{\K}{\mathcal{K}}
\newcommand{\N}{\mathcal{N}}
\renewcommand{\L}{\mathcal{L}}
\providecommand{\Ginf}[1]{\Gamma^{\mathrm{in}}_{#1}}
\title{The back-and-forth relation of the linear order transfers to the labelled tree at a shift of two}
\author{David Gonzalez, Dino Rossegger}
\date{arxiv:2305.07126v1; CC BY 4.0}
\begin{document}
\maketitle
\noindent\textit{Editorial note.} The counted statement is the author environment at source0.tex lines 616-620 and its proof environment opens at line 621, the line immediately after; the pairing is by adjacency and the retained proof range 621-645 is the whole environment.\par\medskip
\noindent\textit{Editorial note.} The retained preamble is the author style file dino.sty (package lines, the theorem declaration support, the notation and tuple macros and the theorem environments) followed by the ordinal and rank macros of the manuscript itself, reproduced verbatim; this is the same recipe as the previously published card of this article, so no package outside it is requested and extra\_packages stays empty.\par\medskip
\noindent\textit{Editorial note.} Metadata provenance: source\_id and source\_version are the values recorded in this source licence evidence file (sources/scott-linear-orders-2305.07126v1/LICENSE\_EVIDENCE.json: arxiv:2305.07126v1 / 2305.07126v1), which the card points at directly; the title and author list are transcribed from the authors\_index and title keys of the source record of the pinned copy. The card carries no counter setting and imposes no original numbering, since this article has no label-number dump and its numbering cannot be reconstructed from the retained ranges alone.\par\medskip
\noindent\textit{Editorial note.} No context span and no external reference are needed: the retained statement and proof use no label or citation from outside their own ranges.\par\medskip
\noindent\textit{Editorial scope.} One original result of the pinned author manuscript, counted once: the transfer of the back-and-forth relation from the linear order of a structure to its labelled tree at a shift of two ranks (author label lem:bf\_treetolo).\par\medskip
\noindent\textit{Allowed standard prerequisites (editorial).}\begin{itemize}
\item The linear order attached to a structure. A structure A carries a linear order L\_A of labelled blocks, and the labelled tree T\_A of its finite tuples is coded inside it.. It is the pair of objects the lemma relates.
\item The coding map f. The map f sends the elements of L\_A to the corresponding nodes of the labelled tree, so that comparing f applied to a tuple compares the coded tuples.. It is the map through which the two relations are compared.
\item Back-and-forth relations of both structures. The rank alpha back-and-forth relations are defined by the game of that rank, and a rank comparison is decided by the interval relations left after each move.. They are the two relations shown to correspond at shifted ranks.
\end{itemize}\medskip
\begin{lemma}\label{lem:bf_treetolo}
   For every $\alpha$ and ordered $\ba,\bb\in L_\A$ such that if $a_i$, or $b_i$ are from a labeled
   block, then every element in this block is in $\ba$, respectively $\bb$, $\ba\leq_{3+\alpha}\bb$ if
   and only if $f(\ba)\leq_{1+\alpha} f(\bb)$.
\end{lemma}

\begin{proof}
   The proof is by transfinite induction with the only interesting case being
   the base case. For this note that the label and predecessor
   relation of $\T_\A$ in $L_{\A}$ are both $\Sinf{3}$ and $\Pinf{3}$
   definable in $L_{\T_\A}$. Let $\phi$ be a $\Sinf{1}$ formula true of
   $\ba$, then this formula can be translated into a $\Sinf{3}$ formula true
   of $f(\ba)$ and thus $\ba\leq_3 \bb$ implies $f(\ba)\leq_1 f(\bb)$. 

   On the other hand, assume that $f(\ba)\leq_1 f(\bb)$ and without loss of
   generality that $\ba$ and $\bb$ contain elements from exactly two
   different blocks. Let us consider the back-and-forth game where the
   $\forall$-player plays elements $\bb^1$ on their first turn. For every
   $b^1_i$ the $\exists$-player looks for elements $a^1_i$ such that
   $l(f(a^1_i))=l(f(b^1_i))$, $a^1_i$ is in the same position in its block as
   $b^1_i$ and $\ba\ba^1$ and $\bb\bb^1$ are
   isomorphic in the tree-ordering. Note that they will find a suitable $\ba^1$
   because $f(\ba)\leq_1 f(\bb)$. Now on their last turn the $\forall$-player
   plays a tuple $\ba^2$ and the $\exists$-player has to respond playing
   $\bb^2$ such that $\ba\ba^1\ba^2\leq_1 \bb\bb^1\bb^2$. All they need to do
   is ensure that the intervals in $\ba\ba^1\ba^2$ are at least as large as the
   intervals in $\bb\bb^1\bb^2$. Between any two elements in different blocks in
   $L_\A$, there are blocks coding paths in $\T_\A$. Thus, between any two
   elements in different blocks there are blocks of arbitrary large size and
   thus the $\exists$-player can find elements satisfying this requirement.
\end{proof}

\end{document}
